\section{Discussion and Limitations}
\label{sec:disc}
\paragraph{What the evidence supports.}
Under a common protocol, target items plus source-model information reduce the error of the full-score estimate by about a fifth to a quarter relative to a random subset, and the correction term makes this robust to extrapolation. It does not show that \mass{} is better than existing methods in general: against the baseline chosen by source-only cross-validation the difference is within noise, anchor-IRT is not distinguishable from \mass{}-Ens in the extrapolation headline table, and under interpolation ridge on a random subset is as good or better. Ranking is not improved, and under extrapolation the error is of the size of the spread among the strongest models.

\paragraph{Limitations.}
(i) \emph{Development choice.} The ensemble predictor of \mass{}-Ens was chosen after a pilot that looked at target-model errors on one split of IFEval, ARC-Challenge, GPQA and GSM; those benchmarks are flagged, the latent-only \mass{} is reported alongside, and the nine non-development benchmarks give a similar summary. The target-side results of the baselines existed before the source-CV baseline selection was computed (the selection itself uses only sources). The method list was frozen before the final sweep; later runs (sensitivity) use a reduced method set and 10 repetitions.
(ii) \emph{One extrapolation split per benchmark}: the strongest 20\% of models; uncertainty comes from model-level bootstrap and item-selection seeds, not from alternative splits. Targets in the top 20\% are few for HELM benchmarks (16--17 models). The family-out split uses the organisation prefix of model names, coarser than a base-model family, and HELM has no family analogue.
(iii) \emph{Dependence}: the six HELM scenarios share models, and benchmark-level intervals treat the benchmarks as independent; we pool no significance statements and report ``wins'' without multiplicity correction. The 13 benchmarks are the unit of the cross-benchmark summaries; bootstrap $B=200$ for model-level intervals.
(iv) \emph{Estimator.} The predictor is fit on the sampled items without cross-fitting, so design-unbiasedness is only approximate; the standard error ignores this and under-covers for $n\le 20$.
(v) \emph{Scope.} Binary scores only; no item text is used; Pacchiardi et al.'s assessor, Active Evaluation Acquisition and the metabench GAM were not run (they need item embeddings, a trained policy or specialised fitting not available here), and \texttt{anchor\_irt}, \texttt{anchor\_aipw} and \texttt{irt\_info} are simplified variants under their own names, so their results do not speak for the original methods. Low-spread benchmarks (GPQA, MuSR, MATH Lvl~5) give poor rankings for every method; we report them instead of dropping them. Source models are open-weight leaderboard models with many near-duplicates and fine-tunes, which makes interpolation easier than evaluating a new architecture.
(vi) \emph{Compute.} Every run is CPU only; the heaviest sweep (MMLU-Pro, 12k items) takes about 15 h of wall time per method group on a loaded 2-CPU host (Appendix~\ref{app:repro}).

\paragraph{Broader impact.}
The work concerns cheaper evaluation of language models; the risk is over-trusting a subset-based score, in particular for ranking or for models unlike the source pool. We recommend reporting the interval and the source pool, and using the full benchmark for decisions about the strongest models.
